% FILE: 01_research_question.tex
% Project: atlas
% Role: doctrine-atlas-and-thesis-chapter-surface-map

\section{Research Question}
\label{sec:atlas-rq}

\subsection{Problem Statement}
\label{sec:atlas-rq-problem}

The process-intelligence research foundry addresses a compound problem: the
absence of a unified, citable doctrine surface that maps the seven distinct
module boundaries of the research program --- Aalst Broadcaster, AGI Adversarial
Defenses, Blue River Dam, GGEN Prompt Manufactory, Supabase Law Substrate, WASM
Boundary Law, and ZoeApp Ministry Ontology --- onto the formal objects of the
Chatman Equation and the lifecycle calculus. Without such a surface, dissertation
chapters would duplicate module definitions, introduce inconsistent naming, and
make cross-module citations unreliable.

The deeper problem the atlas addresses is epistemic: what is the unit of
doctrine in a multi-repository research program? Each implementation repository
(wasm4pm, wasm4pm-compat, ggen, ostar) operates under its own authority; each
has its own receipt chain and its own ALIVE/PARTIAL verdict. But the
dissertation requires a single plane of reference in which all seven doctrine
modules can be cited together, their boundaries respected, and their
contribution to the thesis visible without requiring the reader to navigate seven
separate repositories. The atlas is the answer to that problem.

\subsection{Old-Regime Assumption Challenged}
\label{sec:atlas-rq-old-regime}

The old-regime assumption challenged by the atlas is that documentation is
derivative: a post-hoc description of an implementation that was built first.
Under this assumption, a ``chapter surface map'' is merely a table of contents
assembled after the research is done. It adds no epistemic value; it is a
convenience, not a claim.

The atlas rejects this assumption. Each of the seven doctrine-surface files is
itself a normative artifact: it names a module, bounds its claims, and
establishes the vocabulary through which all downstream citations must pass. The
\texttt{supabase\_law\_substrate.md} file does not merely describe Row Level
Security; it establishes the law that ``React state is not authorization'' as a
citable doctrine point, with specific RLS policy table entries (profiles,
actor\_commands, actor\_events, actor\_receipts, actor\_outbox,
actor\_quarantine) as its evidence base. The atlas-as-map-first position asserts
that naming a doctrine module before its receipt chain is complete is a lawful
act of PARTIAL manufacturing --- not a premature claim to ALIVE.

\subsection{Hypothesis}
\label{sec:atlas-rq-hypothesis}

\textbf{[AUTHOR THESIS]} The hypothesis is that a flat set of seven
doctrine-surface files, each named and bounded at the start of the research
program, provides sufficient navigational and definitional scaffolding to enable
multi-chapter PhD dissertation writing without inter-chapter naming conflicts,
without duplication of module internals, and without requiring the reader to
resolve cross-repository authority disputes. The atlas achieves this by encoding
each module's core claim, its key primitive(s), and its relationship to the
Chatman Equation in a single authoritative file.

The secondary hypothesis is that the PARTIAL status of the atlas is not a
deficiency but a structural property of a doctrine-surface artifact: it maps
modules whose receipts live in other repositories, and its own receipt would be
redundant with --- not additive to --- the PI\_RESEARCH\_PROGRAM\_ALIVE\_001
checkpoint that seals the foundry as a whole.

\subsection{Success Criteria}
\label{sec:atlas-rq-success}

Success for the atlas is defined by three criteria, all of which can be
evaluated without a receipt internal to the atlas directory:

\begin{enumerate}
  \item \textbf{Naming completeness.} All seven doctrine modules are named,
        bounded, and citable. Each file exists and contains at minimum: a module
        name, a core claim, and at least one key primitive. \textit{Status:
        met.} All seven files present and dated 2026-06-01.

  \item \textbf{Foundry ALIVE.} The underlying foundry that the atlas maps
        achieves ALIVE status on all 12 audit gates. \textit{Status: met.}
        PI\_RESEARCH\_PROGRAM\_ALIVE\_001 sealed at 12/12 gates PASS, 342
        artifacts, BLAKE3 \texttt{e7c8f2d94a71b5c3e9f1d6a4b2c8e5f7a1d3c5b7e9f2d4a6c8e0f1a3b5c7d9}.

  \item \textbf{Dissertation usability.} Each atlas file provides a sufficient
        vocabulary for a dissertation chapter to cite without restating module
        internals. \textit{Status: PARTIAL --- the ZoeApp Ministry Ontology and
        AGI Adversarial Defenses files are notably thin and may require
        manufacturing expansion before the corresponding dissertation chapters
        can reach full citation depth.}
\end{enumerate}
